\scanpage{160}\mspage{18}
% archive sheet number: 160
De plus \deleted{si $\unsure{\mathscr X|\mathcal X}$ est loc. connexe,} ces conditions impliquent que $\mathscr X$ est connexe, donc $0$-connexe si $\mathscr X\ne\varnothing$.

\inserted{Cela \uline{implique de plus que} $\mathscr X$ est $\unsure{W_\infty|1W_\infty}$-asph, cf P.S. \S35 (corollaire p. 32)}

% slanted
\margin{$\longleftarrow$ \uuline{Faux} loc. cit. une \unsure{déf.} bien \unsure{différente}}

C'est essentiellement Prop. 1 de P.S. \S65, loc. cit. p. 170--171.$^{*}$ Cf. aussi .PS \S35, \deleted{pages} 50--55,

% slanted
\margin{Mais je m'y \unsure{borne} aux topos de la forme $\mathrm{Top}(A)$, avec $A$ dans $\mathrm{Cat}$.}

\uline{Définition} 2.4. On dit que $\mathscr X$ est \uline{tot. $0$-connexe}, s'il est \deleted{\illegible{2}} loc. connexe, non vide et satisfait les conditions équivalentes de 2.3. (NB c'est un cas particulier de la notion ``totalement $W$-asphérique'' de loc. cit. On a aussi la notion de totalement $n$-connexe.)

Notons que pour une \deleted{catég.} objet $X$ de $\mathrm{Cat}$, dire que $\mathrm{Top}(X)$ est tot. $0$-connexe revient à dire que $X\ne\varnothing$, et que $+$ que $x,y\in\mathrm{Ob}\,X$, $x\times y$ dans $X^{\wedge}$ est $0$-connexe, i.e. $X_{/x\times y}$ est $0$-connexe.$^{*}$ C'est une condition très forte.

\inserted{\unsure{conviendrons|nous dirons} que $X$ est \uline{tot. $0$-connexe}. \unsure{Notons ici}}

% slanted
\margin{Cela signifie aussi que $\operatorname{diag}_X:X\longrightarrow X\times X$ est co-cofinale, i.e. $X^\circ\longrightarrow X^\circ\times X^\circ$ est cofinal, et signifie qu'il en est de $=$ de $X\longrightarrow X^E$ pour tt ens \uline{fini} $E$.}

\uline{Proposition} 2.5. Soit $\mathscr X$ un topos, \deleted{Si $\mathscr X$} Si $\mathscr X$ est tot. $0$-connexe, alors $\mathscr X$ est irréductible. \deleted{\unsure{alors} \illegible{2} \unsure{qu'il}}

(En fait la \unsure{conclusion|condition} \unsure{vaut|est}, \unsure{tout à fait générale|tout à fait gén.}, en supposant seulement que le produit de deux objets $0$-connexes de $\mathscr X^{\wedge}$ est \uline{non} \uline{vide}, et il suffit $=$ de le supposer pour des sous-objets de l'objet final $e_{\mathscr X}$.) En particulier, si $\mathscr X=\mathrm{Top}(X)$, où $X$ est dans $\mathrm{Cat}$, alors $X$ est fbt. filtrante à g.

Mais \deleted{la condit.} bien sûr, \deleted{\unsure{si} $X$ est \unsure{loc.}} cela n'implique \uline{pas} que $\mathscr X$ soit loc. irréductible.
